\documentclass[11pt]{article}

\usepackage[margin=1in]{geometry}
\usepackage{amsmath,amssymb,amsthm}
\usepackage{graphicx}
\usepackage{microtype}
\usepackage[hidelinks]{hyperref}

\newtheorem{theorem}{Theorem}
\newtheorem{lemma}[theorem]{Lemma}
\newcommand{\Ftwo}{\mathbb F_2^+}
\newcommand{\HH}{\mathcal H}
\newcommand{\emptyword}{\varnothing}

\title{The Sharp Free Approximation Rate for $(1-x)(1-y)$\\
\large An exact recursion and a full candidate answer to the example in arXiv:2209.10373}
\author{Automated mathematics run \texttt{fa\_banach\_001}, lane 5}
\date{9 August 2026}

\begin{document}
\maketitle

\begin{center}
\fbox{\parbox{0.91\textwidth}{\textbf{Status.}
Candidate full solution to the explicit sharpness question for
$f(x,y)=(1-x)(1-y)$, likely valid and awaiting human review.  The theorem
below proves the stronger asymptotic $\sqrt n\,c_n\to1$.}}
\end{center}

\section{Source question}

The source is P.~Arora, M.~Augat, M.~T.~Jury, and M.~Sargent,
\emph{An Optimal Approximation Problem for Free Polynomials},
arXiv:2209.10373, published in \emph{Proceedings of the American Mathematical
Society} \textbf{152} (2024), 455--470 \cite{AAJS}.  Free words are declared
orthonormal, and the authors consider
\[
 c_n=\min_{\deg p\le n}\|p f-1\|_2^2.
\]
In Section 3.1 (arXiv PDF page 14), they observe that for
$f(x,y)=(1-x)(1-y)$ one can obtain $c_n\lesssim n^{-1/2}$, but state that
they do not know whether the exponent $1/2$ is sharp.  The relevant source
passage is reproduced below.

\begin{figure}[p]
\centering
\includegraphics[width=\textwidth]{figures/open_problem_crop.png}
\caption{The sharpness question in Section 3.1 of the source, arXiv PDF page 14.}
\end{figure}
\clearpage

\section{Definitions and statement}

Let $\Ftwo$ be the free monoid on the letters $x,y$, with empty word
$\emptyword$.  Write $e_w$ for the word basis of
$\ell^2(\Ftwo)$, and let
\[
 \HH_n=\operatorname{span}\{e_w:|w|\le n\}.
\]
A free polynomial is identified with its coefficient vector in this Hilbert
space.  The optimal squared error for the source example is
\[
 c_n:=\min_{p\in\HH_n}
 \|p(1-x)(1-y)-e_{\emptyword}\|_2^2.
\]

\begin{theorem}\label{thm:main}
Define two positive sequences $(e_m)_{m\ge-1}$ and
$(\kappa_m)_{m\ge-1}$ by
\[
 e_{-1}=1,\qquad \kappa_{-1}=\frac12,
\]
and, for $m\ge0$,
\begin{align}
 \kappa_m
 &=\kappa_{m-1}+\frac{e_{m-1}}{1+e_{m-1}},\label{eq:krec}\\
 e_m
 &=1+\frac{\kappa_{m-1}e_{m-1}}
 {\kappa_{m-1}+e_{m-1}}.\label{eq:erec}
\end{align}
Then the optimal errors for $f(x,y)=(1-x)(1-y)$ satisfy the exact formula
\[
 \boxed{\ c_n=\frac1{e_n}\ }\qquad(n\ge0).
\]
Moreover,
\[
 \boxed{\ \lim_{n\to\infty}\sqrt n\,c_n=1\ }.
\]
Consequently the exponent $p=1/2$ in the source question is sharp, and its
leading constant is one.
\end{theorem}

For orientation, the first exact values are
\[
 c_0=\frac34,\qquad c_1=\frac7{11},\qquad
 c_2=\frac{14}{25},\qquad c_3=\frac{56}{111}.
\]

\section{Proof intuition}

Right multiplication by $(1-x)(1-y)$ sends the coefficient at a word $w$
to a four-point pattern supported at $w,wx,wy,wxy$.  Orthogonality to all
such patterns rooted at depth at most $n$ is therefore the local relation
\[
 h_w-h_{wx}-h_{wy}+h_{wxy}=0.
\]
Projection duality says that $c_n^{-1}$ is the least squared norm of a vector
satisfying all these relations and normalized by $h_{\emptyword}=1$.

This is an electrical-energy-style minimization on a finite binary tree.  A
rooted subtree has only two boundary data relevant to its parent: its root
and its $y$-child.  Eliminating all other values turns the entire tree into a
$2\times2$ positive quadratic form.  Its Schur complement is summarized by
one scalar $\kappa_m$, while the unconstrained-root energy is $e_m$.  This
gives the exact recursion in Theorem~\ref{thm:main}.

For the asymptotics, $e\mapsto1+\kappa e/(\kappa+e)$ has a moving fixed point
$r(\kappa)$ defined by $r^2-r=\kappa$.  The recursion keeps $e_m$ below this
fixed point, while any fixed relative gap is closed by a uniformly positive
amount at each step; the fixed point itself moves only by $O(1/r)$.  Hence
$e_m/r(\kappa_m)\to1$.  Since $\kappa_m\sim m$, the claimed limit follows.

\section{Exact dual recursion}

Let $R_xe_w=e_{wx}$ and $R_ye_w=e_{wy}$ be the right creation operators, and
put
\[
 A=(I-R_y)(I-R_x).
\]
Thus $Ap=p(1-x)(1-y)$ and
\[
 c_n=\operatorname{dist}(e_{\emptyword},A\HH_n)^2.
\]
Let $L_n=(A\HH_n)^\perp$.  Standard Hilbert-space projection duality gives
\begin{equation}\label{eq:dual}
 \frac1{c_n}
 =\min\{\|h\|_2^2:h\in L_n,\ h_{\emptyword}=1\}.
\end{equation}
Indeed, if $g$ is the orthogonal projection of $e_{\emptyword}$ onto $L_n$,
then $c_n=\|g\|^2$; among $h\in L_n$ with
$\langle e_{\emptyword},h\rangle=1$, equality in Cauchy--Schwarz is attained
by $h=g/\|g\|^2$.

For every word $w$,
\[
 (A^*h)_w=h_w-h_{wx}-h_{wy}+h_{wxy}.
\]
Consequently $h\in L_n$ precisely when
\begin{equation}\label{eq:local}
 h_w-h_{wx}-h_{wy}+h_{wxy}=0
 \qquad (|w|\le n).
\end{equation}
Only coordinates of depth at most $n+2$ occur in these constraints.  Setting
all deeper coordinates to zero preserves feasibility and decreases the norm,
so the minimum in \eqref{eq:dual} is a finite-tree problem.

We now solve that problem.  For $m\ge-1$, let $E_m(a)$ be the minimum energy
in a rooted binary subtree whose root has value $a$, with
\eqref{eq:local} imposed through relative depth $m$.  Let $F_m(a,z)$ be the
same minimum when the value $z$ at the root's $y$-child is also fixed.  The
case $m=-1$ means that no relation is imposed, so
\[
 E_{-1}(a)=|a|^2,
 \qquad F_{-1}(a,z)=|a|^2+|z|^2.
\]
Tree symmetry and homogeneity give
\[
 E_m(a)=e_m|a|^2,
 \qquad
 F_m(a,z)=
 \begin{bmatrix}\overline a&\overline z\end{bmatrix}
 Q_m
 \begin{bmatrix}a\\z\end{bmatrix},
\]
where $Q_m$ is a real symmetric positive matrix.  The data are real, so this
Hermitian formulation also shows that allowing complex coefficients does not
change the scalar recursion.

Suppose
\[
 Q_{m-1}=\begin{bmatrix}A&B\\B&C\end{bmatrix},
 \qquad
 \kappa_{m-1}:=
 \frac{\det Q_{m-1}}{A+2B+C}.
\]
At a root with value $a$, write $b,c,d$ for the values at $wx,wy,wxy$.
The root relation is $d=b+c-a$.  If the $y$-child is fixed at $z=c$, the
two child subtrees are disjoint and therefore
\begin{equation}\label{eq:Fsplit}
 F_m(a,z)=|a|^2+E_{m-1}(z)
 +\min_b F_{m-1}(b,b+z-a).
\end{equation}
For any $s\in\mathbb C$, direct completion of the square gives
\[
 \min_b F_{m-1}(b,b+s)
 =\frac{AC-B^2}{A+2B+C}|s|^2
 =\kappa_{m-1}|s|^2.
\]
Hence
\begin{equation}\label{eq:Qrec}
 Q_m=
 \begin{bmatrix}
 1+\kappa_{m-1}&-\kappa_{m-1}\\
 -\kappa_{m-1}&e_{m-1}+\kappa_{m-1}
 \end{bmatrix}.
\end{equation}
If the $y$-child is not fixed, a second one-variable minimization in
\eqref{eq:Fsplit} yields
\[
 E_m(a)
 =|a|^2+\min_z\bigl(e_{m-1}|z|^2
 +\kappa_{m-1}|z-a|^2\bigr)
 =\left(1+\frac{\kappa_{m-1}e_{m-1}}
 {\kappa_{m-1}+e_{m-1}}\right)|a|^2.
\]
This is \eqref{eq:erec}.  From \eqref{eq:Qrec},
\[
 \frac{\det Q_m}{(Q_m)_{11}+2(Q_m)_{12}+(Q_m)_{22}}
 =\kappa_{m-1}+\frac{e_{m-1}}{1+e_{m-1}},
\]
which is \eqref{eq:krec}.  Since $Q_{-1}=I$, its initial Schur-complement
parameter is $\kappa_{-1}=1/2$.

Finally, the minimum in \eqref{eq:dual} is exactly the root energy
$E_n(1)=e_n$.  Therefore $c_n=1/e_n$, proving the exact part of
Theorem~\ref{thm:main}.

\section{Sharp asymptotics}

We prove that $e_n\sim\sqrt n$.  Write $k_n=\kappa_n$ for brevity and define
\[
 r(k)=\frac{1+\sqrt{1+4k}}2,
 \qquad r(k)^2-r(k)=k.
\]
For fixed $k>0$, set
\[
 \Phi_k(u)=1+\frac{ku}{k+u}.
\]
The number $r(k)$ is the positive fixed point of $\Phi_k$, and $\Phi_k$ is
increasing.  Since $e_{-1}=1<r(k_{-1})$ and $k_n>k_{n-1}$, induction gives
\begin{equation}\label{eq:belowroot}
 1\le e_n\le r(k_n)\qquad(n\ge-1).
\end{equation}
Indeed, $e\le r(k)$ implies
$\Phi_k(e)\le r(k)<r(k')$ for every $k'>k$.  It also implies monotonicity,
because
\begin{equation}\label{eq:increment}
 \Phi_k(e)-e=\frac{k+e-e^2}{k+e}\ge0.
\end{equation}

The recursion for $k_n$, together with $e_n\ge1$, gives
\begin{equation}\label{eq:kbounds}
 \frac{n+2}{2}\le k_n<n+\frac32
 \qquad(n\ge0).
\end{equation}
We also record a useful elementary lower bound.  Since
$k_0=1$ and $e_0=4/3$, induction shows
\begin{equation}\label{eq:halfsqrt}
 e_n\ge\frac12\sqrt{k_n}\qquad(n\ge0).
\end{equation}
To verify the induction step, suppose $k\ge1$ and
$e\ge\tfrac12\sqrt k$, and put
$k'=k+e/(1+e)\le k+1$.  If
$e\ge\tfrac12\sqrt{k+1}$ there is nothing to prove, by
\eqref{eq:increment}.  Otherwise $e^2<(k+1)/4$ and
$e<\tfrac12\sqrt{k+1}\le k$, so
\[
 \Phi_k(e)-e
 =\frac{k+e-e^2}{k+e}
 >\frac{k-(k+1)/4}{2k}\ge\frac14.
\]
But
\[
 \frac12(\sqrt{k+1}-\sqrt k)
 =\frac1{2(\sqrt{k+1}+\sqrt k)}<\frac14,
\]
which proves $\Phi_k(e)\ge\tfrac12\sqrt{k+1}\ge
\tfrac12\sqrt{k'}$.

Equations \eqref{eq:belowroot}--\eqref{eq:halfsqrt} already give explicit
two-sided estimates
\begin{equation}\label{eq:explicit}
 \frac1{1+\sqrt{n+3/2}}
 \le c_n\le\frac{2\sqrt2}{\sqrt{n+2}}.
\end{equation}
Thus the exponent $1/2$ is sharp even without identifying the leading
constant.  We finish by proving that the constant is one.

The lower bounds above imply $e_n\to\infty$.  Hence
\[
 k_n=\frac12+\sum_{j=-1}^{n-1}\frac{e_j}{1+e_j}
 \quad\Longrightarrow\quad \frac{k_n}{n}\longrightarrow1
\]
by Ces\`aro convergence.  It remains to compare $e_n$ with $r(k_n)$.

Fix $0<\varepsilon<1$ and set $s=(1-\varepsilon)r$, where $r=r(k)$.
The increment $\Phi_k(u)-u$ is decreasing in $u>0$, and at $s$ it equals
\begin{equation}\label{eq:boundaryincrement}
 \Phi_k(s)-s
 =\frac{(r-s)(r+s-1)}{k+s}
 =\varepsilon\frac{(2-\varepsilon)r-1}{r-\varepsilon}
 \longrightarrow\varepsilon(2-\varepsilon)>0.
\end{equation}
On the other hand, since $k_{n+1}-k_n<1$ and
$r'(k)=1/(2r(k)-1)$,
\begin{equation}\label{eq:rootmotion}
 r(k_{n+1})-r(k_n)
 \le\frac1{2r(k_n)-1}\longrightarrow0.
\end{equation}
Choose $d>0$ and $N$ so that, for $n\ge N$, the boundary increment in
\eqref{eq:boundaryincrement} is at least $d$, while the increase of
$(1-\varepsilon)r(k_n)$ is less than $d/2$.  Whenever
$e_n\le(1-\varepsilon)r(k_n)$, its next increment is at least $d$, so its
gap below that moving boundary decreases by at least $d/2$.  It crosses the
boundary after finitely many steps.  Once it has crossed, monotonicity of
$\Phi_{k_n}$ and the same two estimates show that it remains above the
boundary.  Therefore
\[
 \liminf_{n\to\infty}\frac{e_n}{r(k_n)}\ge1-\varepsilon.
\]
Since \eqref{eq:belowroot} gives the reverse upper bound and $\varepsilon$ is
arbitrary,
\[
 \frac{e_n}{r(k_n)}\longrightarrow1.
\]
Now $r(k_n)/\sqrt{k_n}\to1$ and $k_n/n\to1$, so
$e_n/\sqrt n\to1$.  The identity $c_n=1/e_n$ completes the proof of
Theorem~\ref{thm:main}.

\section{Verification notes}

The exact recursion was checked independently against the original primal
least-squares problem.  The script \texttt{code/verify\_recurrence.py} builds
the four-term coefficient column of $w(1-x)(1-y)$ for every $|w|\le n$,
forms its Gram matrix over the rational numbers, solves the normal equations,
and compares the resulting residual with $1/e_n$.  The values agree exactly
for every tested degree.  The script also iterates the scalar recurrence in
floating point to display convergence of $\sqrt n/e_n$ to one.  These checks
are corroborative only; no numerical assertion is used in the proof.

The most delicate formal points are local and exposed in the packet: the
right-multiplication order in the definition of $A$, the fact that $wxy$ is
the $y$-child of the $x$-child in the tree split, and the Schur-complement
coefficient in the one-variable minimization.  The separate
\texttt{verification.md} records these checks and the tested range.

\section{Scope and novelty check}

This result fully answers the source paper's explicit question whether the
exponent $1/2$ is sharp for $(1-x)(1-y)$.  It does not answer the broader
question in the same section asking whether there is a uniform positive
exponent valid independently of the free polynomial $F$.

The bounded novelty search was performed on 9 August 2026.  The run's
registry, solution, attempt, and proof-gap indexes were searched for
arXiv:2209.10373, the source title, ``optimal polynomial approximants'',
``free polynomial'', and the explicit factor $(1-x)(1-y)$.  Web/arXiv
searches used the exact title and example, the author names, the journal DOI
10.1090/proc/16474, and combinations with ``sharp'', ``rate'', and
``one-half''.  The available citation trail surfaced later work on classical
Hardy-space optimal approximants that cites the source as background, but no
paper treating this free example or giving the recursion above.  The search
is bounded rather than exhaustive, so bibliographic novelty remains subject
to specialist review.

\section{Human-review recommendation}

Review as a full solution to the concrete source example.  First check
\eqref{eq:dual} and the coordinate identity for $A^*$.  Then check the
subtree split \eqref{eq:Fsplit}, particularly that the fixed value $b+z-a$
is the $y$-child of the $x$-child.  Finally verify the moving-boundary
argument after \eqref{eq:boundaryincrement}.  If those points pass, the
exact recurrence and the asymptotic constant follow without any external
mathematical input.

\begin{thebibliography}{9}

\bibitem{AAJS}
P. Arora, M. Augat, M. T. Jury, and M. Sargent,
\emph{An Optimal Approximation Problem for Free Polynomials},
Proc. Amer. Math. Soc. \textbf{152} (2024), no.~2, 455--470;
arXiv:2209.10373,
\href{https://doi.org/10.1090/proc/16474}{doi:10.1090/proc/16474}.

\end{thebibliography}

\end{document}
